\SourceLeaf{296}
\begin{center}
\textbf{NACHLASS UND BRIEFWECHSEL.}

\vspace{1em}

\textbf{[ZUR SPHÄRISCHEN TRIGONOMETRIE.]}

\vspace{0.5em}

\end{center}

\vspace{1em}

\begin{center}[1.]\end{center}

\begin{center}
\textsc{Gauss} an \textsc{Gerling}. Göttingen, 18. Februar 1815.
\end{center}

Ich danke Ihnen sehr, lieber \textsc{Gerling}, für die Mitteilungen aus
Ihrer Arbeit über die Trigonometrie. Ihren Beweis für die 4 Gleichungen finde
ich für Ihren Zweck sehr passend. Zur Vollständigkeit fehlt ihm nur,
dass die Vorzeichen noch gesondert bewiesen werden müssen,
da zum\,Beispiel aus
\[
\sin \tfrac12 A^2 \cos \tfrac12 (b+c)^2 = \cos \tfrac12 (B+C)^2 \cos \tfrac12 a^2
\]
eigentlich nur folgt, dass
\[
\begin{array}{r@{\ }c@{\ }l}
\sin \tfrac12 A \cos \tfrac12 (b+c) & \text{entweder} & {}=\cos \tfrac12 (B+C) \cos \tfrac12 a \\[4pt]
 & \text{oder} & {}=-\cos \tfrac12 (B+C) \cos \tfrac12 a.
\end{array}
\]

Eine solche Ergänzung des Beweises kann man aus anderen Sätzen herleiten,
doch wird dadurch das Ganze weitschweifiger. In letzter Zeit habe ich mich immer an
die folgende Herleitungsmethode gehalten.

Ich setze die 5 Gleichungen voraus, da sie tatsächlich ohnehin alle
notwendig sind.
